\section{The harmonic quotient metric and unitary rigidity}
\label{sec:harmonic-quotient}

We now vary the holomorphic bundle, while keeping the smooth
base fibre fixed. The resulting harmonic metric has a different
domain from the metric in \eqref{eq:finite-flow-metric}.
For a skew-Hermitian connection one-form use
$\|a\|_2^2=-\int\tr(a\wedge *a)$, with volume $\omega_p^3/3!$.
On the matrix space $W=\operatorname{Mat}_2\oplus\operatorname{Mat}_2$
put
\begin{equation}\label{eq:flat-arrow-metric}
 G_{0,p}(q,q)=2c_+(p)\|q_A\|_{\rm HS}^2
                                  +2c_-(p)\|q_B\|_{\rm HS}^2.
\end{equation}
Let $\mathfrak k_0=\mathrm i\mathfrak h$ and let $Q$ also
denote the split complementary projection on skew-Hermitian
trace-free sections. The fundamental vector for the compact
action is
\[
 \rho_w\xi=(\xi_+A-A\xi_-,\ \xi_-B-B\xi_+).
\]
Exact equivariance of the actual connection map gives
\begin{equation}\label{eq:actual-gauge-action}
       d\mathcal A_p(w)(\rho_w\xi)=-D_{\mathcal A_p(w)}\xi.
\end{equation}
The ineffective scalar may equivalently be included in the
finite gauge algebra.

\begin{theorem}\label{thm:harmonic-quotient}
Uniformly in the base and every small $w$, for all $q\in W$
and $v\in QH^1(\mathfrak{su}E^{\rm sp})$,
\begin{equation}\label{eq:all-gauge-competitors}
 \begin{split}
 (1-C|w|)\{G_{0,p}(q,q)+\|D_0v\|_2^2\}
 &\le\|d\mathcal A_p(w)q-D_{\mathcal A_p(w)}v\|_2^2\\
 &\le(1+C|w|)\{G_{0,p}(q,q)+\|D_0v\|_2^2\}.
 \end{split}
\end{equation}
Consequently the full unitary-gauge minimum is relatively
equivalent, with error $O(|w|)$, to
\begin{equation}\label{eq:finite-quotient-metric}
             g^{\rm quot}_{0,p,w}(q)
                    =\inf_{\xi\in\mathfrak k_0}G_{0,p}(q-\rho_w\xi,q-\rho_w\xi).
\end{equation}
Its kernel is exactly $\im\rho_w$.
For $F_p(z)=\mathcal A_p(\chi_p(z))$ and genuine tangents
to a smooth orbit-type stratum of $\mu_{2,p}^{-1}(0)$,
the actual harmonic connection metric satisfies
\begin{equation}\label{eq:harmonic-relative-metric}
 (1-C|z|)g^{\rm quot}_{0,p,z}
       \le F_p^*g_{\rm harm}\le(1+C|z|)g^{\rm quot}_{0,p,z}.
\end{equation}
This real metric is twice the Dolbeault harmonic
Kodaira--Spencer metric. The estimates allow every change
of stabilizer and all positive singular-value ratios of
$\rho_z$.
\end{theorem}

\begin{proof}
Split harmonicity and type orthogonality give
\[
 \|a_{1,p}(q)\|_2^2=2\|\eta_p(q)\|_2^2=G_{0,p}(q,q),
 \qquad \langle a_{1,p}(q),D_0v\rangle_2=0.
\]
The connection expansion \eqref{eq:connection-expansion}
and the split section gap imply
\[
 \begin{split}
 \|d\mathcal A_p(w)q-D_{\mathcal A_p(w)}v
                           -(a_{1,p}(q)-D_0v)\|_2
 &\le\|d\mathcal N_p(w)q\|_2+\|[a_p(w),v]\|_2\\
 &\le C|w|\{G_{0,p}(q,q)+\|D_0v\|_2^2\}^{1/2}.
 \end{split}
\]
Squaring proves \eqref{eq:all-gauge-competitors}; density
extends it to the stated $H^1$ domain. The minimum over
$QH^1$ exists uniquely by coercivity. Infimizing the lower
bound and taking $v=0$ for the upper bound gives a relative
comparison with $G_{0,p}(q,q)$.

Every full gauge direction is $\xi+v$. By
\eqref{eq:actual-gauge-action}, its residual is
$d\mathcal A_p(w)(q+\rho_w\xi)-D_{\mathcal A_p(w)}v$.
Apply the preceding comparison to every $q+\rho_w\xi$
and then infimize in $\xi$. Changing its sign gives
\eqref{eq:finite-quotient-metric}. The finite range is
closed at each parameter, so the kernel is precisely
$\im\rho_w$. No estimate for the size of a minimizing
$\xi$ has been used.

Real equivariance gives
$d\chi_p(z)\rho_z\xi=\rho_{\chi_p(z)}\xi$.
Since $d\chi_p=I+O(|z|)$ in the uniformly positive arrow
metric, the same comparison transfers through $\chi_p$
inside every residual before minimization. This proves
\eqref{eq:harmonic-relative-metric} as a quotient form.

For its harmonic interpretation, take a smooth curve in
one zero-level orbit-type stratum. Its actual derivative
$a$ at the HE connection $D_A$ satisfies
$(D_Aa)^{0,2}=0$ and $\mathrm i\Lambda D_Aa=0$.
The full gauge minimum supplies the Coulomb representative
$a_h=a-D_A\zeta$, with $D_A^*a_h=0$.
Subtracting this vertical preserves the two linearized
equations because $F_A^{0,2}=0$ and
$\mathrm i\Lambda F_A=0$. The Kähler identities identify
the Coulomb and contracted equations with the sum and
difference of the two adjoint expressions, giving
\[
 \dbar_A(a_h^{0,1})=0,\qquad \dbar_A^*(a_h^{0,1})=0.
\]
This is the harmonic representative of the actual
Kodaira--Spencer class; a unitary vertical is a Dolbeault
coboundary. Type orthogonality gives the factor two in
its real norm, in agreement with
\cite[Section 7.2 and equations (7.6.5), (7.6.20)--(7.6.21)]{Kobayashi}.
The orbit-type pieces of the zero level and their quotients
by $K$ are smooth by \cite[Theorem 2.1]{SL}.
The form is positive on their quotient tangents by its
identified kernel. For fixed $p$ it is smooth on each
stratum, by the elliptic inverse on the complement of
the constant-dimensional parallel kernel. No uniform
derivative across strata or through the singular base
is asserted.
\end{proof}

The same action occurs in the weak endomorphism energy:
\begin{equation}\label{eq:shared-action-energy}
             E_{p,z}(S)=G_{0,p}(\rho_z(\mathrm iS),\rho_z(\mathrm iS)).
\end{equation}
Thus its kernel is the infinitesimal stabilizer, while
the null directions of the quotient tangent metric are
the image of the action. These are different spaces
related by one specified map.

\subsection{All unitary intertwiners are constant}

An infinitesimal metric comparison does not identify
different finite quotient points with different bundle
classes. The following statement supplies that identification.

\begin{theorem}\label{thm:unitary-rigidity}
On a common small parameter ball, every unitary gauge
transformation intertwining $\mathcal A_p(w_1)$ and
$\mathcal A_p(w_2)$ is a split parallel block unitary
$K\in U(2)\times U(2)$, and $w_2=K\cdot w_1$.
At HE points, two constructed holomorphic bundles are
isomorphic if and only if their parameters lie in the
same $K$-orbit. The same assertion holds in the balanced
$z$ coordinates.
\end{theorem}

\begin{proof}
Put $\rho=\max(|w_1|,|w_2|)$ and abbreviate the connection
increments by $a_j=a_p(w_j)$. The intertwining equation gives
\[
 D_0k=ka_1-a_2k,\qquad \|D_0k\|_\infty\le C\rho.
\]
Equation~\eqref{eq:connection-divergence} holds for every
small slice parameter: its proof did not use the HE equation.
Taking divergence therefore gives
\[
 L_0k=kD_0^*a_1-(D_0^*a_2)k
       -\sum_j(D_{0,e_j}k)a_1(e_j)
       +\sum_ja_2(e_j)D_{0,e_j}k=O_\infty(\rho^2).
\]
Let $\widehat\Pi$ project onto the full complex split
parallel algebra $\operatorname{Mat}_2\oplus\operatorname{Mat}_2$.
The full complex inverse in Lemma~\ref{lem:graph-inverse}
implies
\begin{equation}\label{eq:intertwiner-complement}
                    \|k-\widehat\Pi k\|_{C^1(D_0)}\le C\rho^2.
\end{equation}
Since $k$ is unitary, $B=\widehat\Pi k$ satisfies
$B^*B=I+O(\rho^2)$. Its block polar factor $K_0$
is unitary and $\|K_0^{-1}k-I\|_{C^1(D_0)}\le C\rho^2$.
In the finite skew-Hermitian block algebra solve
\[
                 \widehat\Pi\log(e^{-C}K_0^{-1}k)=0.
\]
The derivative in $C$ at the identity is $-I$ and differs
from it by $O(\rho^2)$ for these data. A uniform finite
contraction gives $|C|\le C\rho^2$. Thus
\begin{equation}\label{eq:intertwiner-factorization}
 k=Ke^\xi,\qquad K=K_0e^C\in K,\qquad
 \widehat\Pi\xi=0,\qquad \|\xi\|_{C^1(D_0)}\le C\rho^2.
\end{equation}
The determinant connection is fixed, so $\det k$ is
constant. The small logarithm makes $\tr\xi$ constant,
and its zero parallel projection makes this trace zero.

Conjugate the second connection by $K^{-1}$ and put
$\delta=K^{-1}\cdot w_2-w_1$. The segment between these
two parameters stays in the small slice, giving
\[
 a(K^{-1}\cdot w_2)-a(w_1)=a_{1,p}(\delta)+R_\delta,
                         \qquad \|R_\delta\|_2\le C\rho|\delta|.
\]
The exact exponential gauge identity is
\[
 e^\xi D(w_1)e^{-\xi}-D(w_1)
                 =-\int_0^1e^{t\xi}(D(w_1)\xi)e^{-t\xi}\,\dd t.
\]
The split gap on $\widehat\Pi^\perp$, together with
\eqref{eq:intertwiner-factorization}, therefore gives
\[
 \|a_{1,p}(\delta)+D_0\xi\|_2
                         \le C\rho(|\delta|+\|D_0\xi\|_2).
\]
The two terms on the left are orthogonal, and their
squared norm is
$G_{0,p}(\delta,\delta)+\|D_0\xi\|_2^2$.
The uniform positive arrow masses allow absorption for
small $\rho$. Hence $\delta=0$, $D_0\xi=0$, and then
$\xi=0$. Every intertwiner is the constant $K$.

For a holomorphic isomorphism $T$ between two HE members,
the induced Hom connection has zero mean curvature.
Its section Kähler identity gives
$\|D_{\rm Hom}T\|_2^2=2\|\dbar_{\rm Hom}T\|_2^2=0$;
see also \cite[Theorem 3.1.9(i)]{Kobayashi}.
Thus $T$ and its positive polar factor are parallel,
and its unitary polar factor is an intertwiner.
The first assertion applies, without simplicity or
stability assumptions. The converse is equivariance.
Finally injectivity and equivariance of $\chi_p$ transfer
the assertion to balanced parameters.
\end{proof}

\subsection{The topology of the actual local locus}

On a fixed smooth fibre take the affine space of continuous
unitary connections, with its $C^0$ topology, modulo $C^1$
unitary gauge. This quotient is Hausdorff. Indeed, suppose
$D_{j,n}\to D_j$ and $D_{2,n}=k_nD_{1,n}k_n^{-1}$.
The intertwiners are pointwise unitary and
\[
 D_0k_n=k_n(D_{1,n}-D_0)-(D_{2,n}-D_0)k_n.
\]
Their first derivatives are bounded on this fixed compact
manifold. Arzelà--Ascoli gives a uniformly convergent
subsequence; the equation then gives convergence of the
first derivatives and a $C^1$ unitary intertwiner in the
limit. Hence the orbit relation is closed. The orbit
projection is open because saturations are unions of
translates, so the quotient is Hausdorff. For smooth
endpoints the displayed equation bootstraps an intertwiner
to smoothness, introducing no extra identifications.

For any compact invariant subset of the balanced zero
set, Theorem~\ref{thm:unitary-rigidity} makes its quotient
map into actual connection classes continuous and
injective. A compact-to-Hausdorff injection is a
homeomorphism onto its image. These points are precisely
distinct actual holomorphic bundle classes.
The local completeness proved in Section~\ref{sec:balancing}
shows that a nested such locus contains representatives
of every sufficiently small polystable object in the
specified local germ. This is the germ coverage already
proved there; $\chi_p$ is a real parametrization and
does not supply a holomorphic chart or an algebraic
Cartesian square of moduli spaces.
